% ── ملخص (العربية) ────────────────────────────────────────────
\chapter*{{\arabicfont ملخص}}
\addcontentsline{toc}{chapter}{\texorpdfstring{{\arabicfont ملخص}}{ملخص}}

\begin{arabicblock}
يتناول هذا المشروع الختامي، المنجز داخل شركة \textbf{CIVCO BTP}،
تصميم وتطوير تطبيق ويب مخصص لإدارة وتخطيط مشاريع البناء الخاصة بالشركة.
وبسبب تزايد عدد الورش وتعدد الأطراف المعنية،
سعت \textbf{CIVCO BTP} إلى امتلاك أداة رقمية مركزية
تُنظّم النشاط اليومي وتُسهّل تبادل المعلومات بين مختلف الأقسام.

\vspace{0.35cm}

يغطي التطبيق دورة حياة المشروع بالكامل: إنشاء الورش ومتابعتها حسب المراحل،
إدارة الزبناء وجهات الاتصال، تعيين الفرق، تخطيط المهام،
إعداد العروض والفواتير، إنتاج بونات التسليم،
متابعة العقود، والاطلاع على لوحة معلومات موحدة.
يهدف إلى استبدال الأدوات المتفرقة (جداول البيانات، الوثائق الورقية،
التواصل غير الرسمي) بمنصة واحدة
يمكن للأدوار المهنية المختلفة في الشركة الوصول إليها
وفق صلاحياتهم.

\vspace{0.35cm}

تعتمد الحل على بنية \textit{full-stack}:
خادم API مطوَّر بـ \textbf{Laravel} و\textbf{PHP}
يُوفّر خدمات \textbf{REST} مؤمّنة
ويعتمد على قاعدة بيانات علائقية،
مقترنًا بواجهة مستخدم حديثة بـ \textbf{React}
منظمة كـ \textbf{SPA}.
يضمن نظام الأدوار والصلاحيات (\textbf{RBAC})
أن يصل كل مستخدم فقط إلى الوظائف المرتبطة بملفه.

\vspace{0.35cm}

يقدم هذا التقرير سياق المشروع، والخيارات التقنية،
والتصميم الوظيفي والتقني للنظام،
إضافة إلى تنفيذ الوحدات الرئيسية والتحقق منها.
وينقسم إلى أربعة فصول: الإطار العام والمنهجية،
حالة الفن، التصميم \textbf{UML} للنظام،
ثم تنفيذ الوظائف المسلَّمة.

\vspace{0.35cm}

تشمل الوحدات التشغيلية المسلَّمة إدارة المشاريع حسب المراحل
مع خرائط \textbf{Leaflet}،
إنتاج ومتابعة العروض والفواتير وبونات التسليم،
بوابة الزبناء لاستشارة الوثائق والتحقق منها،
المراسلة الداخلية، تخطيط المهام،
متابعة العقود والملحقات،
إضافة إلى وحدة ضبط الأدوار والصلاحيات.
كما يضمن سجل النشاط تتبع العمليات الحساسة داخل التطبيق.

\vspace{0.5cm}
\noindent\textbf{الكلمات المفتاحية:}
البناء، إدارة المشاريع، تطبيق ويب، ERP، Laravel، React،
قطاع البناء، العروض، الفوترة، التخطيط.
\end{arabicblock}
